\section{Конструктивные числа}

\begin{definition}
    Число $x \in \R$ называется \emph{конструктивным}, если на числовой прямой его можно получить из единичного отрезка конечной последовательностью
    построений циркулем и линейкой.
    Число $x \in \R$ называется \emph{конструктивным}, если на числовой прямой его можно получить из единичного отрезка конечной последовательностью
    построений циркулем и линейкой.

    В каждом шаге разрешено строить прямую через две уже построенные точки, окружность с уже построенным центром и уже построенным радиусом,
    а также брать новые точки как пересечения уже построенных прямых и окружностей.
    В каждом шаге разрешено строить прямую через две уже построенные точки, окружность с уже построенным центром и уже построенным радиусом,
    а также брать новые точки как пересечения уже построенных прямых и окружностей.
\end{definition}

\begin{remark}
    Алгебраическая идея построений состоит в том, что новые координаты появляются только из пересечений прямых и окружностей.
    Пересечение прямых задаётся линейной системой, а пересечение с окружностью сводится к квадратному уравнению.
    Поэтому один геометрический шаг добавляет не больше одного квадратного корня.
    Алгебраическая идея построений состоит в том, что новые координаты появляются только из пересечений прямых и окружностей.
    Пересечение прямых задаётся линейной системой, а пересечение с окружностью сводится к квадратному уравнению.
    Поэтому один геометрический шаг добавляет не больше одного квадратного корня.
\end{remark}

\begin{definition}
    Пусть $F \subset \R$ --- подполе. Точка называется $F$-точкой, если обе её координаты лежат в $F$.

    $F$-прямая --- это прямая, проходящая через две различные $F$-точки. Эквивалентно, она задаётся уравнением
    \[
        ax + by + c = 0, \qquad a,b,c \in F,
    \]
    где $(a,b) \neq (0,0)$.

    $F$-окружность --- это окружность с центром в $F^2$ и радиусом из $F$. Эквивалентно, она задаётся уравнением
    \[
        (x-a)^2 + (y-b)^2 = r^2, \qquad a,b,r \in F.
    \]
\end{definition}

\begin{stmt}
    Если прямая проходит через две точки $(x_1,y_1)$ и $(x_2,y_2)$ из $F^2$, то её уравнение можно выбрать с коэффициентами из $F$.

    \begin{proof}
        Такая прямая задаётся уравнением
        \[
            (y_2-y_1)x + (x_1-x_2)y + (x_2y_1-x_1y_2)=0.
        \]
        Все коэффициенты лежат в $F$, так как $F$ замкнуто относительно сложения, вычитания и умножения.
    \end{proof}
\end{stmt}

\begin{stmt}
    Пусть $F \subset \R$ --- поле. Если точка $P=(x,y)$ получается как пересечение двух $F$-прямых, $F$-прямой и
    $F$-окружности или двух $F$-окружностей, то существует $d \in F$ такое, что
    \[
        x,y \in F(\sqrt d).
    \]

    \begin{proof}
        Если пересекаются две прямые, то координаты точки пересечения находятся из линейной системы с коэффициентами из
        $F$. Если точка единственная, определитель системы ненулевой, и по формулам Крамера координаты лежат в $F$.
        Если пересекаются две прямые, то координаты точки пересечения находятся из линейной системы с коэффициентами из
        $F$. Если точка единственная, определитель системы ненулевой, и по формулам Крамера координаты лежат в $F$.

        Если пересекаются прямая и окружность, то одну координату можно линейно выразить через другую и подставить в уравнение окружности.
        Получается квадратное уравнение над $F$. Его корни лежат в поле вида $F(\sqrt D)$, где $D \in F$ --- дискриминант.

        Если пересекаются две окружности,
        \[
            (x-a)^2+(y-b)^2=r^2, \qquad (x-a')^2+(y-b')^2=(r')^2,
        \]
        то после вычитания уравнений получается линейное уравнение над $F$. Значит, этот случай сводится к пересечению прямой и окружности.
    \end{proof}
\end{stmt}

\begin{stmt}
    Если число $\alpha \in \R$ конструктивно, то существует башня полей
    \[
        \Q = F_0 \subset F_1 \subset \cdots \subset F_n \subset \R,
        \qquad F_{i+1}=F_i(\sqrt{d_i}), \quad d_i \in F_i,
    \]
    такая, что $\alpha \in F_n$.

    \begin{proof}
        В начале построения доступны рациональные точки, поэтому стартовым полем координат является $\Q$.
        Каждый новый геометрический шаг, по предыдущему утверждению, добавляет координаты, лежащие либо в текущем поле,
        либо в его квадратичном расширении. После конечного числа шагов все построенные координаты лежат в башне квадратичных расширений над $\Q$.
    \end{proof}
\end{stmt}

\begin{cor}
    Если $\alpha \in \R$ конструктивно и алгебраично над $\Q$, то степень $[\Q(\alpha):\Q]$ делит некоторую степень двойки.

    \begin{proof}
        По предыдущему утверждению $\alpha \in F_n$, где $F_n$ получается из $\Q$ башней квадратичных расширений. Поэтому
        \[
            [F_n:\Q] = [F_n:F_{n-1}]\cdots [F_1:F_0]
        \]
        является степенью двойки или делителем степени двойки. Так как $\Q(\alpha) \subset F_n$, по мультипликативности степеней
        $[\Q(\alpha):\Q]$ делит $[F_n:\Q]$.
        является степенью двойки или делителем степени двойки. Так как $\Q(\alpha) \subset F_n$, по мультипликативности степеней
        $[\Q(\alpha):\Q]$ делит $[F_n:\Q]$.
    \end{proof}
\end{cor}

\section{Классические невозможности}

\begin{example}[Удвоение куба]
    Чтобы построить куб объёма $2$ из единичного куба, нужно построить ребро длины $\sqrt[3]{2}$.

    Минимальный многочлен числа $\sqrt[3]{2}$ над $\Q$ равен
    \[
        x^3 - 2.
    \]
    Он неприводим над $\Q$ по критерию Эйзенштейна при $p=2$. Его степень равна $3$, а степень конструктивного алгебраического числа должна делить
    степень двойки.
    Следовательно, $\sqrt[3]{2}$ не конструктивно, и удвоение куба невозможно.
\end{example}

\begin{example}[Трисекция угла $60^\circ$]
    Если бы угол $60^\circ$ можно было разделить на три равные части, то был бы конструктивен угол $20^\circ$,
    а значит, было бы конструктивно число $\cos 20^\circ$.

    Положим $x=2\cos20^\circ$. Из формулы
    \[
        \cos(3\varphi)=4\cos^3\varphi-3\cos\varphi
    \]
    при $\varphi=20^\circ$ получаем
    \[
        4\cos^3 20^\circ - 3\cos20^\circ = \frac12.
    \]
    После замены $x=2\cos20^\circ$ это даёт
    \[
        x^3-3x-1=0.
    \]
    Многочлен $x^3-3x-1$ не имеет рациональных корней, поэтому как кубический многочлен он неприводим над
    $\Q$. Значит, $[\Q(x):\Q]=3$, что невозможно для конструктивного числа. Следовательно, трисекция угла
    $60^\circ$ невозможна.
    Многочлен $x^3-3x-1$ не имеет рациональных корней, поэтому как кубический многочлен он неприводим над
    $\Q$. Значит, $[\Q(x):\Q]=3$, что невозможно для конструктивного числа. Следовательно, трисекция угла
    $60^\circ$ невозможна.
\end{example}

\begin{example}[Квадратура круга]
    Чтобы построить квадрат площади $\pi$, нужно построить сторону длины $\sqrt\pi$. Если бы это было возможно,
    то $\sqrt\pi$ было бы конструктивным, а значит, алгебраическим. Тогда и $\pi=(\sqrt\pi)^2$ было бы алгебраическим.
    Но $\pi$ трансцендентно. Следовательно, квадратура круга невозможна.
    Чтобы построить квадрат площади $\pi$, нужно построить сторону длины $\sqrt\pi$. Если бы это было возможно,
    то $\sqrt\pi$ было бы конструктивным, а значит, алгебраическим. Тогда и $\pi=(\sqrt\pi)^2$ было бы алгебраическим.
    Но $\pi$ трансцендентно. Следовательно, квадратура круга невозможна.
\end{example}

\section{Критерий Эйзенштейна}

\begin{stmt}
    Пусть
    \[
        f(x)=a_nx^n+a_{n-1}x^{n-1}+\cdots+a_1x+a_0 \in \Z[x].
    \]
    Если существует простое число $p$ такое, что
    \[
        p \nmid a_n, \qquad p\mid a_i \text{ при } 0\le i<n, \qquad p^2\nmid a_0,
    \]
    то $f$ неприводим над $\Q$.

    \begin{proof}
        По лемме Гаусса достаточно доказать неприводимость в $\Z[x]$ для примитивного многочлена. Предположим противное:
        \[
            f(x)=g(x)h(x), \qquad g,h\in \Z[x], \qquad \deg g,\deg h>0.
        \]
        Переходя по модулю $p$, получаем
        \[
            \overline f(x)=\overline a_n x^n \in \F_p[x], \qquad \overline a_n\neq0.
        \]
        Значит, $\overline g$ и $\overline h$ являются мономами с ненулевыми старшими коэффициентами.
        Поэтому свободные члены $g(0)$ и $h(0)$ делятся на $p$. Тогда $a_0=f(0)=g(0)h(0)$ делится на
        $p^2$, что противоречит условию.
        Значит, $\overline g$ и $\overline h$ являются мономами с ненулевыми старшими коэффициентами.
        Поэтому свободные члены $g(0)$ и $h(0)$ делятся на $p$. Тогда $a_0=f(0)=g(0)h(0)$ делится на
        $p^2$, что противоречит условию.
    \end{proof}
\end{stmt}

\section{Корни из единицы и циклотомические многочлены}

\begin{definition}
    Число
    \[
        \zeta_p = e^{2\pi i/p}
    \]
    называется примитивным $p$-м корнем из единицы, если $\zeta_p^p=1$ и $\zeta_p^k\neq1$ при $1\le k<p$.
\end{definition}

\begin{stmt}
    Для простого $p$ многочлен
    \[
        \Phi_p(x)=x^{p-1}+x^{p-2}+\cdots+x+1=\frac{x^p-1}{x-1}
    \]
    неприводим над $\Q$.

    \begin{proof}
        Рассмотрим замену $x\mapsto x+1$:
        \[
            \Phi_p(x+1)=\frac{(x+1)^p-1}{x}
            = \sum_{k=1}^{p}\binom pk x^{k-1}
            = x^{p-1}+p x^{p-2}+\cdots+\binom p2x+p.
        \]
        Все коэффициенты, кроме старшего, делятся на $p$, а свободный член равен $p$ и не делится на
        $p^2$. По критерию Эйзенштейна $\Phi_p(x+1)$ неприводим над $\Q$. Следовательно, и $\Phi_p(x)$ неприводим над $\Q$.
    \end{proof}
\end{stmt}

\begin{stmt}
    Если $p$ --- простое число, то
    \[
        [\Q(\zeta_p):\Q]=p-1.
    \]

    \begin{proof}
        Примитивный корень $\zeta_p$ является корнем многочлена $\Phi_p(x)$. Так как $\Phi_p(x)$ неприводим над
        $\Q$, он является минимальным многочленом $\zeta_p$ над $\Q$. Поэтому степень расширения равна
        $\deg \Phi_p=p-1$.
        Примитивный корень $\zeta_p$ является корнем многочлена $\Phi_p(x)$. Так как $\Phi_p(x)$ неприводим над
        $\Q$, он является минимальным многочленом $\zeta_p$ над $\Q$. Поэтому степень расширения равна
        $\deg \Phi_p=p-1$.
    \end{proof}
\end{stmt}

\begin{stmt}
    Пусть $p$ --- нечётное простое число. Тогда
    \[
        \left[\Q\left(\cos\frac{2\pi}{p}\right):\Q\right]=\frac{p-1}{2}.
    \]

    \begin{proof}
        Положим
        \[
            \alpha=\zeta_p+\zeta_p^{-1}=2\cos\frac{2\pi}{p}.
        \]
        Тогда $\alpha\in\Q(\zeta_p)$, а $\zeta_p$ является корнем квадратного уравнения
        \[
            t^2-\alpha t+1=0
        \]
        над $\Q(\alpha)$. Значит, $[\Q(\zeta_p):\Q(\alpha)]\le2$.

        Так как $p$ нечётно, $\zeta_p\notin\R$, но $\alpha\in\R$. Следовательно, $\Q(\zeta_p)\neq\Q(\alpha)$,
        и степень $[\Q(\zeta_p):\Q(\alpha)]$ равна $2$. Поэтому
        \[
            [\Q(\alpha):\Q]=\frac{[\Q(\zeta_p):\Q]}{[\Q(\zeta_p):\Q(\alpha)]}=\frac{p-1}{2}.
        \]
        Умножение на $2$ не меняет поле, значит, та же степень получается для $\cos(2\pi/p)$.
    \end{proof}
\end{stmt}

\section{Правильные многоугольники}

\begin{definition}
    Правильный $n$-угольник построим циркулем и линейкой тогда и только тогда, когда можно построить центральный угол
    $2\pi/n$. Алгебраически это сводится к построимости числа
    Правильный $n$-угольник построим циркулем и линейкой тогда и только тогда, когда можно построить центральный угол
    $2\pi/n$. Алгебраически это сводится к построимости числа
    \[
        \cos\frac{2\pi}{n}.
    \]
\end{definition}

\begin{stmt}
    Если правильный $p$-угольник построим для нечётного простого $p$, то
    \[
        \frac{p-1}{2}
    \]
    является степенью двойки.

    \begin{proof}
        Если правильный $p$-угольник построим, то число $\cos(2\pi/p)$ конструктивно. Его степень над
        $\Q$ равна $(p-1)/2$. По необходимому условию конструктивности эта степень должна делить степень двойки.
        Следовательно, $(p-1)/2$ само является степенью двойки.
        Если правильный $p$-угольник построим, то число $\cos(2\pi/p)$ конструктивно. Его степень над
        $\Q$ равна $(p-1)/2$. По необходимому условию конструктивности эта степень должна делить степень двойки.
        Следовательно, $(p-1)/2$ само является степенью двойки.
    \end{proof}
\end{stmt}

\begin{stmt}
    Если $p$ --- простое число и $p=2^m+1$, то $m$ является степенью двойки.

    \begin{proof}
        Пусть $m=ab$, где $b$ нечётно и $b>1$. Тогда
        \[
            2^m+1=(2^a)^b+1.
        \]
        При нечётном $b$ выражение $u^b+1$ делится на $u+1$. Значит, $2^m+1$ делится на $2^a+1$ и потому составно.
        Противоречие с простотой $p$. Следовательно, у $m$ нет нечётных делителей больше $1$, то есть
        $m$ является степенью двойки.
        При нечётном $b$ выражение $u^b+1$ делится на $u+1$. Значит, $2^m+1$ делится на $2^a+1$ и потому составно.
        Противоречие с простотой $p$. Следовательно, у $m$ нет нечётных делителей больше $1$, то есть
        $m$ является степенью двойки.
    \end{proof}
\end{stmt}

\begin{definition}
    Простые числа вида
    \[
        F_k=2^{2^k}+1
    \]
    называются простыми числами Ферма.
\end{definition}

\begin{example}[Пятиугольник]
    Для $p=5$ имеем
    \[
        \left[\Q\left(\cos\frac{2\pi}{5}\right):\Q\right]=2.
    \]
    Действительно,
    \[
        \cos\frac{2\pi}{5}=\frac{\sqrt5-1}{4}.
    \]
    Поэтому правильный пятиугольник построим.
\end{example}

\begin{example}[Семнадцатиугольник]
    Число $17=2^{2^2}+1$ является простым числом Ферма. Поэтому правильный $17$-угольник построим. Нужный косинус выражается через вложенные
    квадратные радикалы,
    что и соответствует построимости циркулем и линейкой.
\end{example}

\begin{stmt}
    Правильный $n$-угольник построим циркулем и линейкой тогда и только тогда, когда
    \[
        n = 2^a p_1p_2\cdots p_s,
    \]
    где $a\ge0$, а $p_1,\ldots,p_s$ --- попарно различные простые числа Ферма.
\end{stmt}

\begin{remark}
    На лекции для простого $p$ было разобрано необходимое условие через степень поля $\Q(\cos(2\pi/p))$.
    Обратное направление удобно доказывать через теорию Галуа: расширения степени $2^n$ раскладываются в башни квадратичных расширений.
\end{remark}

\section{$p$-группы}

\begin{definition}
    Конечная группа $G$ называется $p$-группой, если
    \[
        |G|=p^n
    \]
    для некоторого простого $p$ и некоторого $n\ge0$.
\end{definition}

\begin{definition}
    Центр группы $G$ --- это подгруппа
    \[
        Z(G)=\{z\in G \mid zg=gz \text{ для всех } g\in G\}.
    \]
    Централизатор элемента $x\in G$ --- это подгруппа
    \[
        C_G(x)=\{g\in G\mid gx=xg\}.
    \]
\end{definition}

\begin{definition}
    Группа $G$ действует на себе сопряжениями:
    \[
        g\cdot x = gxg^{-1}.
    \]
    Орбита элемента $x$ при этом действии называется классом сопряжённости:
    \[
        \Orb(x)=\{gxg^{-1}\mid g\in G\}.
    \]
    Стабилизатор элемента $x$ при этом действии равен $C_G(x)$.
\end{definition}

\begin{stmt}
    Для любого $x\in G$ верно
    \[
        |\Orb(x)|=[G:C_G(x)].
    \]

    \begin{proof}
        Это формула орбита--стабилизатор для действия сопряжением. Стабилизатором элемента $x$ является множество элементов,
        коммутирующих с $x$, то есть $C_G(x)$.
        Это формула орбита--стабилизатор для действия сопряжением. Стабилизатором элемента $x$ является множество элементов,
        коммутирующих с $x$, то есть $C_G(x)$.
    \end{proof}
\end{stmt}

\begin{stmt}
    Для конечной группы $G$ верно уравнение классов
    \[
        |G| = |Z(G)| + \sum_i [G:C_G(x_i)],
    \]
    где $x_i$ пробегают представителей нецентральных классов сопряжённости.

    \begin{proof}
        Элементы центра --- это ровно элементы, классы сопряжённости которых состоят из одного элемента.
        Остальные элементы разбиваются на классы сопряжённости размеров $[G:C_G(x_i)]$. Складывая размеры всех классов,
        получаем нужную формулу.
        Элементы центра --- это ровно элементы, классы сопряжённости которых состоят из одного элемента.
        Остальные элементы разбиваются на классы сопряжённости размеров $[G:C_G(x_i)]$. Складывая размеры всех классов,
        получаем нужную формулу.
    \end{proof}
\end{stmt}

\begin{stmt}
    Если $G$ --- конечная $p$-группа и $|G|>1$, то центр $Z(G)$ нетривиален.

    \begin{proof}
        В уравнении классов
        \[
            |G| = |Z(G)| + \sum_i [G:C_G(x_i)]
        \]
        левая часть делится на $p$. Если $x_i\notin Z(G)$, то $C_G(x_i)\neq G$, поэтому индекс $[G:C_G(x_i)]$ больше
        $1$. Так как $G$ --- $p$-группа, такой индекс является степенью $p$ и потому делится на $p$.

        Значит, все слагаемые в сумме делятся на $p$. Тогда и $|Z(G)|$ делится на $p$. В частности, центр содержит не только единичный элемент.
    \end{proof}
\end{stmt}

\section{Разрешимые группы}

\begin{definition}
    Группа $G$ называется разрешимой, если существует цепочка подгрупп
    \[
        G=G_0 \triangleright G_1 \triangleright \cdots \triangleright G_n=\{e\}
    \]
    такая, что $G_{i+1}$ нормальна в $G_i$, а все фактор-группы $G_i/G_{i+1}$ абелевы.
\end{definition}

\begin{example}[Абелева группа]
    Если $G$ абелева, то она разрешима: достаточно взять цепочку
    \[
        G\triangleright \{e\}.
    \]
    Фактор $G/\{e\}\cong G$ абелев.
\end{example}

\begin{stmt}
    Если $N\subset Z(G)$ и группа $G/N$ разрешима, то $G$ разрешима.

    \begin{proof}
        Возьмём разрешимый ряд для $G/N$ и поднимем все его члены в $G$ по естественной проекции $G\to G/N$.
        Факторы в поднятом ряду изоморфны факторам исходного ряда, поэтому они абелевы. В конце остаётся цепочка
        $N\triangleright\{e\}$, и фактор $N$ абелев, так как $N\subset Z(G)$.
    \end{proof}
\end{stmt}

\begin{stmt}
    Любая конечная $p$-группа разрешима.

    \begin{proof}
        Докажем индукцией по $|G|$. При $|G|=1$ утверждение очевидно. Пусть $|G|>1$. У конечной $p$-группы центр нетривиален,
        поэтому в $Z(G)$ есть подгруппа $N$ порядка $p$. Такая подгруппа центральна, а значит, нормальна.
        Докажем индукцией по $|G|$. При $|G|=1$ утверждение очевидно. Пусть $|G|>1$. У конечной $p$-группы центр нетривиален,
        поэтому в $Z(G)$ есть подгруппа $N$ порядка $p$. Такая подгруппа центральна, а значит, нормальна.

        Фактор $G/N$ снова является $p$-группой и имеет меньший порядок. По индукции $G/N$ разрешима.
        Так как $N\subset Z(G)$, из предыдущего утверждения следует, что $G$ разрешима.
        Фактор $G/N$ снова является $p$-группой и имеет меньший порядок. По индукции $G/N$ разрешима.
        Так как $N\subset Z(G)$, из предыдущего утверждения следует, что $G$ разрешима.
    \end{proof}
\end{stmt}

\begin{remark}
    Эта групповая часть нужна для обратного направления в задачах о построимости. Башня квадратичных расширений соответствует цепочке подгрупп индекса
    $2$ в группе Галуа. Поэтому свойства $2$-групп напрямую связаны с построениями циркулем и линейкой.
    Эта групповая часть нужна для обратного направления в задачах о построимости. Башня квадратичных расширений соответствует цепочке подгрупп индекса
    $2$ в группе Галуа. Поэтому свойства $2$-групп напрямую связаны с построениями циркулем и линейкой.
\end{remark}

\section{Квадратичные расширения}

\begin{stmt}
    Пусть $E/F$ --- расширение полей степени $2$ и $\Char F\neq2$. Тогда существует $d\in F$ такое, что
    \[
        E=F(\sqrt d).
    \]

    \begin{proof}
        Возьмём элемент $\gamma\in E\setminus F$. Тогда $E=F(\gamma)$, потому что $[E:F]=2$.

        Минимальный многочлен $\gamma$ над $F$ имеет степень $2$:
        \[
            x^2+bx+c, \qquad b,c\in F.
        \]
        Его корни выражаются через дискриминант:
        \[
            \gamma=\frac{-b\pm\sqrt{b^2-4c}}{2}.
        \]
        Так как $\Char F\neq2$, число $2$ обратимо в $F$. Значит,
        \[
            E=F(\gamma)=F\bigl(\sqrt{b^2-4c}\bigr).
        \]
    \end{proof}
\end{stmt}

\begin{stmt}
    Пусть
    \[
        \Q = F_0 \subset F_1 \subset \cdots \subset F_n \subset \R,
        \qquad F_{i+1}=F_i(\sqrt{d_i}).
    \]
    Тогда каждый элемент поля $F_n$ является конструктивным числом.

    \begin{proof}
        Над уже построенным полем можно геометрически выполнять сложение, вычитание, умножение, деление и извлечение квадратного корня из положительного элемента.
        Эти операции реализуются стандартными построениями с подобными треугольниками и окружностями.
        Индукция по $i$ показывает, что все элементы каждого поля $F_i$ конструктивны.
        Над уже построенным полем можно геометрически выполнять сложение, вычитание, умножение, деление и извлечение квадратного корня из положительного элемента.
        Эти операции реализуются стандартными построениями с подобными треугольниками и окружностями.
        Индукция по $i$ показывает, что все элементы каждого поля $F_i$ конструктивны.
    \end{proof}
\end{stmt}

\section{Галуа-расширения степени степени двойки}

\begin{stmt}
    Пусть $E\subset\R$, расширение $E/\Q$ является конечным галуа-расширением и
    \[
        [E:\Q]=2^n.
    \]
    Тогда все элементы $E$ конструктивны.

    \begin{proof}
        Группа Галуа
        \[
            G=\Gal(E/\Q)
        \]
        имеет порядок $2^n$, то есть является $2$-группой. У конечной $2$-группы можно построить цепочку подгрупп
        \[
            G=G_0 \triangleright G_1 \triangleright \cdots \triangleright G_n=\{e\},
            \qquad [G_i:G_{i+1}]=2.
        \]
        По соответствию Галуа ей отвечает башня промежуточных полей
        \[
            \Q = E^{G_0} \subset E^{G_1} \subset \cdots \subset E^{G_n}=E,
        \]
        причём каждая следующая степень равна $2$:
        \[
            [E^{G_{i+1}}:E^{G_i}]=2.
        \]
        Каждое расширение степени $2$ получается добавлением квадратного корня. Значит, $E$ лежит в башне квадратичных расширений над
        $\Q$, и все элементы $E$ конструктивны.
        Каждое расширение степени $2$ получается добавлением квадратного корня. Значит, $E$ лежит в башне квадратичных расширений над
        $\Q$, и все элементы $E$ конструктивны.
    \end{proof}
\end{stmt}

\begin{stmt}
    Если $p=2^{2^k}+1$ --- простое число Ферма, то правильный $p$-угольник построим циркулем и линейкой.

    \begin{proof}
        Пусть $\zeta_p=e^{2\pi i/p}$. Группа Галуа циклотомического поля имеет вид
        \[
            \Gal(\Q(\zeta_p)/\Q)\cong(\Z/p\Z)^*.
        \]
        Она циклическая порядка $p-1=2^{2^k}$. Вещественное подполе
        \[
            \Q(\zeta_p+\zeta_p^{-1})=\Q\left(2\cos\frac{2\pi}{p}\right)
        \]
        имеет степень
        \[
            \frac{p-1}{2}=2^{2^k-1}.
        \]
        Это галуа-расширение над $\Q$ степени степени двойки. По предыдущему утверждению все его элементы конструктивны,
        в частности $\cos(2\pi/p)$. Следовательно, правильный $p$-угольник построим.
        Это галуа-расширение над $\Q$ степени степени двойки. По предыдущему утверждению все его элементы конструктивны,
        в частности $\cos(2\pi/p)$. Следовательно, правильный $p$-угольник построим.
    \end{proof}
\end{stmt}

\begin{example}[Почему $7$-угольник не строится]
    Для $p=7$ имеем
    \[
        \left[\Q\left(\cos\frac{2\pi}{7}\right):\Q\right]=3.
    \]
    Степень $3$ не является степенью двойки, поэтому правильный $7$-угольник не построим циркулем и линейкой.
\end{example}

\begin{example}[Почему $9$-угольник не строится]
    Число $9$ не имеет вид $2^a$ умножить на произведение различных простых Ферма: простое $3$ входит в него с квадратом.
    Поэтому правильный $9$-угольник не построим. Эквивалентно, для построения $9$-угольника пришлось бы выполнить трисекцию угла
    $60^\circ$.
    Число $9$ не имеет вид $2^a$ умножить на произведение различных простых Ферма: простое $3$ входит в него с квадратом.
    Поэтому правильный $9$-угольник не построим. Эквивалентно, для построения $9$-угольника пришлось бы выполнить трисекцию угла
    $60^\circ$.
\end{example}
